% =============================================================================
%  §6  定性图景（总结）+ 术语表与记号 + 参考文献 + 习题解答（附录）
%  指挥官撰写。校验器锚点：Glossary / Notation / References / Bibliographical notes / Exercises。
% =============================================================================
\TLsection{6 \ 定性图景与总结}

% ── 五章合一（一）────────────────────────────────────────────────────────────
\begin{frame}{定性图景：五章合一（一）——一条主线}
  Villani 第 I 部分的"定性描述"沿一条主线展开：从\emph{解的存在}到\emph{解的几何与正则性}。
  \begin{center}
  \resizebox{\linewidth}{!}{%
  \begin{tikzpicture}[>={Stealth[length=6pt]},node distance=6mm,
      b/.style={draw=TLprimary,fill=TLprimaryTint,rounded corners=3pt,align=center,
        inner sep=5pt,text width=2.5cm,font=\scriptsize}]
    \node[b] (e) {\textbf{存在}\\§1 · Ch.4\\最优耦合存在};
    \node[b,right=of e] (d) {\textbf{刻画}\\§2 · Ch.5\\对偶 + 循环单调};
    \node[b,right=of d] (g) {\textbf{几何}\\§3 · Ch.6\\$\Wp{p}$ 距离};
    \node[b,right=of g] (k) {\textbf{动力学}\\§4 · Ch.7\\位移插值 = 测地线};
    \node[b,right=of k] (r) {\textbf{正则性}\\§5 · Ch.8\\缩短 $\Rightarrow$ 内部光滑};
    \draw[->,TLaccent,thick] (e)--(d);
    \draw[->,TLaccent,thick] (d)--(g);
    \draw[->,TLaccent,thick] (g)--(k);
    \draw[->,TLaccent,thick] (k)--(r);
  \end{tikzpicture}}
  \end{center}
  \begin{itemize}\small
    \item \textbf{存在}（§1）：紧性 + 下半连续 $\Rightarrow$ 最优耦合存在。
    \item \textbf{刻画}（§2）：最优 $\iff$ c-循环单调 $\iff$ 集中在 $\partial_c\psi$ $\iff$ 对偶等号。
    \item \textbf{几何}（§3）：$\Wp{p}$ 把最优成本升格为距离，$\Ppp{p}(\X)$ 是 Polish 空间。
    \item \textbf{动力学}（§4）：位移插值 = $\Wp{p}$ 测地线。\textbf{正则性}（§5）：轨迹不相交 $\Rightarrow$ 内部绝对连续。
  \end{itemize}
  \TLtakeaway{五章是一条因果链：存在 $\to$ 对偶刻画 $\to$ 度量几何 $\to$ 测地线动力学 $\to$ 内部正则性。}
\end{frame}

% ── 五章合一（二）────────────────────────────────────────────────────────────
\begin{frame}{定性图景：五章合一（二）——核心公式一览}
  \begin{center}\small
  \begin{tabular}{ll}
    \toprule
    \textbf{主题} & \textbf{核心命题 / 公式} \\
    \midrule
    Kantorovich 问题 & $C(\mu,\nu)=\inf_{\pi\in\coup{\mu}{\nu}}\int c\dd\pi$ \\
    对偶（Thm 5.10） & $C(\mu,\nu)=\sup_{\psi\text{ c-凸}}\bigl[\int\ctr{\psi}\dd\nu-\int\psi\dd\mu\bigr]$ \\
    最优性刻画 & $\pi$ 最优 $\iff$ $\supp\pi$ c-循环单调 $\iff$ $\pi$ 集中在 $\partial_c\psi$ \\
    Brenier（二次成本） & $T=\nabla\varphi$，$\varphi$ 凸（$\mu$ 绝对连续） \\
    Wasserstein 距离 & $\Wp{p}(\mu,\nu)=\bigl(\inf_\pi\int d^p\dd\pi\bigr)^{1/p}$ \\
    度量化（Thm 6.9） & $\Wp{p}(\mu_k,\mu)\to0 \iff \mu_k\wkto\mu$ 且 $p$ 阶矩收敛 \\
    位移插值（Thm 7.21） & $\mu_t=(e_t)_\#\Pi$，$\Wp{p}(\mu_s,\mu_t)=|t-s|\,\Wp{p}(\mu_0,\mu_1)$ \\
    缩短 / 正则（Thm 8.7） & $\mu_0$ 绝对连续 $\Rightarrow$ $\mu_t$ 绝对连续（$0<t<1$） \\
    \bottomrule
  \end{tabular}
  \end{center}
  \TLtakeaway{这张表是全讲义的"一页纸总览"：把它读懂，即掌握了最优传输定性理论的骨架。}
\end{frame}

% ── 术语表 · Glossary ─────────────────────────────────────────────────────────
\begin{frame}{术语表 · Glossary（一）}
  \begin{center}\small
  \begin{tabular}{ll}
    \toprule
    \textbf{术语（Term）} & \textbf{含义} \\
    \midrule
    耦合 / 转移计划 (coupling) & 以 $\mu,\nu$ 为边际的 $\X\times\Y$ 上概率测度 $\pi$ \\
    推前 (pushforward) $\pf{T}{\mu}$ & $T$ 把质量从 $\mu$ 搬运后得到的测度 \\
    Monge / Kantorovich 问题 & 确定性映射 / 可分裂耦合下的最优传输 \\
    c-循环单调 (c-cyclically monotone) & 无法通过有限循环重路由降低成本的支撑 \\
    c-凸 / c-变换 (c-convex / c-transform) & Legendre 变换的推广；$\ctr{\psi}(y)=\inf_x[\psi(x)+c(x,y)]$ \\
    Kantorovich 势 (potential) & 对偶问题的最优 c-凸函数 $\psi$ \\
    c-次微分 $\partial_c\psi$ & $\{(x,y):\ctr{\psi}(y)-\psi(x)=c(x,y)\}$，最优匹配关系 \\
    Brenier 映射 & 二次成本下的最优映射 $T=\nabla\varphi$（$\varphi$ 凸） \\
    \bottomrule
  \end{tabular}
  \end{center}
  \TLtakeaway{§1--§2 的语言：耦合、推前、对偶、c-凸与 c-次微分——理解它们即理解最优传输的"语法"。}
\end{frame}

\begin{frame}{术语表 · Glossary（二）}
  \begin{center}\small
  \begin{tabular}{ll}
    \toprule
    \textbf{术语（Term）} & \textbf{含义} \\
    \midrule
    Wasserstein 距离 $\Wp{p}$ & 以 $d^p$ 为成本的最优传输值的 $p$ 次根 \\
    Wasserstein 空间 $\Ppp{p}(\X)$ & $p$ 阶矩有限的概率测度，配以 $\Wp{p}$（Polish） \\
    Kantorovich–Rubinstein & $\Wp{1}=\sup_{\norm{\psi}_{\mathrm{Lip}}\le1}(\int\psi\dd\mu-\int\psi\dd\nu)$ \\
    胶着 (tight) & 质量不逃逸到无穷（Prokhorov 紧性判据） \\
    作用量 / Lagrange 量 (action) & $A(\gamma)=\int_0^1 L(\gamma_t,\dot\gamma_t,t)\dd t$ \\
    常速测地线 (constant-speed geodesic) & $d(\gamma_s,\gamma_t)=|t-s|\,d(\gamma_0,\gamma_1)$ \\
    位移插值 (displacement interpolation) & $\mu_t=(e_t)_\#\Pi$，$\Wp{p}$ 空间中的测地线 \\
    缩短原理 (shortening principle) & 最优轨迹不相交 $\Rightarrow$ 内部正则性 \\
    \bottomrule
  \end{tabular}
  \end{center}
  \TLtakeaway{§3--§5 的语言：距离、测地线、作用量与缩短——把静态的传输升格为动态的几何。}
\end{frame}

\begin{frame}{术语表 · Glossary（三）——§0 预备概念}
  \begin{center}\small
  \begin{tabular}{ll}
    \toprule
    \textbf{术语（Term）} & \textbf{含义} \\
    \midrule
    下半连续 (lsc) & $\liminf_{x\to x_0}f(x)\ge f(x_0)$；等价于所有下水平集 $\{f\le a\}$ 为闭集 \\
    上半连续 (usc) & $\limsup_{x\to x_0}f(x)\le f(x_0)$；lsc 的对偶概念 \\
    绝对连续 $\Pac$ & $\mu\ll\mathrm{vol}$（有 Lebesgue 密度）；无原子；Brenier 定理的前提 \\
    Hausdorff 测度 $\mathcal{H}^{k}$ & 度量空间上的 $k$ 维体积；$\mathcal{H}^{1}=$ 弧长，$\mathcal{H}^{n}\propto\mathrm{Leb}^{n}$ \\
    条件测度 / disintegration & 耦合 $\pi$ 沿边际 $\mu$ 分解为 $\{\pi_{x}\}_{x\in\X}$；Gluing Lemma 的基础 \\
    \bottomrule
  \end{tabular}
  \end{center}
  \TLtakeaway{§0 背景术语：lsc 保证最优性问题有解；$\Pac$ 保证 Brenier 映射存在；disintegration 支撑 Gluing Lemma。}
\end{frame}

% ── 记号表 · Notation ────────────────────────────────────────────────────────
\begin{frame}{记号表 · Notation}
  \begin{center}\small
  \begin{tabular}{ll}
    \toprule
    \textbf{记号（Symbol）} & \textbf{含义} \\
    \midrule
    $\X,\Y$ & Polish（完备可分度量）空间 \\
    $\Pp(\X)$，$\Ppp{p}(\X)$ & 概率测度全体；$p$ 阶矩有限者 \\
    $\Pac(\X)$ & 绝对连续概率测度（有 Lebesgue 密度）\\
    $\coup{\mu}{\nu}=\Pi(\mu,\nu)$ & $(\mu,\nu)$ 的全体耦合（转移计划） \\
    $\pf{T}{\mu}$ & 推前测度（$T$ 把 $\mu$ 搬成的测度） \\
    $c(x,y)$，$C(\mu,\nu)$ & 单位成本；最优总成本 \\
    $\ctr{\psi}$，$\cctr{\psi}$，$\partial_c\psi$ & c-变换、二次 c-变换、c-次微分 \\
    $\Wp{p}$，$\Wtwo$ & $p$ 阶、$2$ 阶 Wasserstein 距离 \\
    $\mu_k\wkto\mu$ & 弱（窄）收敛 \\
    $e_t$，$\Pi$，$\mu_t$ & 取值映射 $\gamma\mapsto\gamma_t$；路径测度；位移插值 \\
    \bottomrule
  \end{tabular}
  \end{center}
  \TLtakeaway{所有记号在 §0 预备一节给出严格定义；本表供随时回查。}
\end{frame}

% ── 参考文献 · References ─────────────────────────────────────────────────────
\begin{frame}{参考文献 · References（一）：来源与原始文献}
  \textbf{主要来源（source）。}
  \begin{itemize}\small
    \item C. Villani, \emph{Optimal Transport: Old and New}, Grundlehren der math.\ Wissenschaften \textbf{338}, Springer, 2009. ——本讲义覆盖第 I 部分（第 4--8 章）。
  \end{itemize}
  \textbf{原始文献与经典（primary literature）。}
  \begin{itemize}\small
    \item G. Monge, \emph{Mémoire sur la théorie des déblais et des remblais}, 1781.
    \item L.V. Kantorovich, \emph{On the translocation of masses}, Dokl.\ Akad.\ Nauk SSSR, 1942.
    \item R.T. Rockafellar, \emph{Characterization of subdifferentials of convex functions}, 1966.
    \item Y. Brenier, \emph{Polar factorization and monotone rearrangement of vector-valued functions}, CPAM, 1991.
    \item R.J. McCann, \emph{A convexity principle for interacting gases}, Adv.\ Math., 1997.
    \item J.-D. Benamou, Y. Brenier, \emph{A computational fluid mechanics solution to the Monge--Kantorovich problem}, Numer.\ Math., 2000.
  \end{itemize}
\end{frame}

\begin{frame}{参考文献 · References（二）：教材与综述}
  \textbf{延伸阅读（further reading / textbooks）。}
  \begin{itemize}\small
    \item C. Villani, \emph{Topics in Optimal Transportation}, AMS, 2003.
    \item L. Ambrosio, N. Gigli, G. Savaré, \emph{Gradient Flows in Metric Spaces and in the Space of Probability Measures}, Birkhäuser, 2008.
    \item F. Santambrogio, \emph{Optimal Transport for Applied Mathematicians}, Birkhäuser, 2015.
    \item G. Peyré, M. Cuturi, \emph{Computational Optimal Transport}, Found.\ Trends ML, 2019.
    \item A. Figalli, F. Glaudo, \emph{An Invitation to Optimal Transport, Wasserstein Distances, and Gradient Flows}, EMS, 2021.
  \end{itemize}
  \TLtakeaway{从 Villani 的两本书入门定性理论，AGS 深入梯度流，Santambrogio / Peyré--Cuturi 面向应用与计算。}
\end{frame}

% ── 文献注记 · Bibliographical notes ──────────────────────────────────────────
\begin{frame}{文献注记 · Bibliographical Notes（按章）}
  \begin{itemize}\small
    \item \textbf{第 4 章（存在性）。} 最优耦合的存在性是经典结果；关键工具 Prokhorov 紧性定理见 Billingsley。
    \item \textbf{第 5 章（对偶）。} 对偶源于 Kantorovich (1942)；c-循环单调与 c-凸推广了 Rockafellar (1966) 的凸分析；势函数构造借鉴 Rüschendorf；二次成本的 Brenier (1991)、Knott--Smith 同期独立给出。
    \item \textbf{第 6 章（Wasserstein）。} 距离 $\Wp{1}$ 即 Kantorovich--Rubinstein；"Wasserstein" 之名经 Dobrushin 等推广（实由 Vasershtein 引入）。
    \item \textbf{第 7 章（位移插值）。} 位移插值与位移凸性由 McCann (1997) 提出；动力学（连续性方程）表述见 Benamou--Brenier (2000)，Otto 据此发展 $\Ppp{2}$ 上的形式 Riemann 几何。
    \item \textbf{第 8 章（缩短原理）。} 不相交思想可追溯至 Monge；定量"缩短引理"源于 Mather (1991) 的作用量极小测度理论，由 Bernard--Buffoni 引入最优传输。
  \end{itemize}
  \TLtakeaway{这套理论由 Monge（1781）发端，经 Kantorovich（1942）线性化，到 Brenier/McCann（1990s）几何化，再到 Otto/Villani（2000s）联结 Ricci 曲率——本讲义止于其"定性"地基。}
\end{frame}

% =============================================================================
\appendix

\begin{frame}{附录 · Appendix：习题解答 · Exercise Solutions}
  本附录收录 §0--§5 全部习题的完整解答。建议先独立完成、再对照。
  \TLtakeaway{解答是用来\emph{核对}的，不是用来\emph{替代}思考的——先挣扎，后对照。}
\end{frame}

% ── §0 解答 ──────────────────────────────────────────────────────────────────
\begin{frame}{习题解答 · §0（预备）}
  \textbf{0.1}~(\diffI)\ $T(x)=2x$，$\mu=\mathrm{Unif}[0,1]$。对 $y\in(0,2)$，$\nu[(0,y)]=\mu[2x<y]=\mu[x<y/2]=y/2$，故密度 $g(y)=\tfrac12$ 于 $[0,2]$，即 $\nu=\mathrm{Unif}[0,2]$。

  \medskip
  \textbf{0.2}~(\diffII)\ 乘积测度满足 $(\mu\otimes\nu)[A\times\Y]=\mu[A]\nu[\Y]=\mu[A]$，同理第二边际为 $\nu$，故 $\mu\otimes\nu\in\coup{\mu}{\nu}$。于是 $\coup{\mu}{\nu}\neq\varnothing$（Kantorovich 可行集总非空）。

  \medskip
  \textbf{0.3}~(\diffIII)\ $\mu=\delta_0$ 的质量必须分裂到 $\pm1$，无函数 $T$ 能做到，故 Monge 不可行。唯一耦合 $\pi=\tfrac12\delta_{(0,-1)}+\tfrac12\delta_{(0,1)}$，$C(\mu,\nu)=\tfrac12\cdot1+\tfrac12\cdot1=1$。

  \TLtakeaway{§0：推前是换元；乘积测度保证可行集非空；点质量必须分裂时 Monge 失效而 Kantorovich 仍解。}
\end{frame}

% ── §1 解答 ──────────────────────────────────────────────────────────────────
\begin{frame}{习题解答 · §1（基本性质）}
  \textbf{1.1}~(\diffI)\ $\mu\otimes\nu$ 的边际为 $\mu,\nu$（见 0.2）。由 Fubini（$c\ge0$）：$I[\mu\otimes\nu]=\iint c\dd\mu\dd\nu<+\infty$。故 $C(\mu,\nu)\le I[\mu\otimes\nu]<+\infty$。

  \medskip
  \textbf{1.2}~(\diffII)\ 设子计划 $\tilde\pi=\pi'/\pi'[A\times B]$ 不最优，则有同边际而更便宜的 $\hat\pi$。令 $Z=\pi'[A\times B]$，构造 $\pi_{\mathrm{new}}=(\pi^\star-\pi')+Z\hat\pi$：边际与 $\pi^\star$ 相同，成本 $=I[\pi^\star]-Z\,I[\tilde\pi]+Z\,I[\hat\pi]<I[\pi^\star]$，与 $\pi^\star$ 最优矛盾。

  \medskip
  \textbf{1.3}~(\diffIII)\ 例 4.9 中唯一最优 Kantorovich 计划把每个 $(0,a)$ 的质量对半分到 $(\pm1,a)$，不是任何映射的图，故无 Monge 极小子；但用"几乎处处单侧"的确定性映射逼近，其成本趋于 $C(\mu,\nu)$，故 $\inf_{\text{Monge}}=\min_{\text{Kantorovich}}$。

  \TLtakeaway{§1：独立耦合给出有限上界；限制性质靠"局部替换"的反证；Monge 下确界可等于 Kantorovich 最小值却不被取到。}
\end{frame}

% ── §2 解答 ──────────────────────────────────────────────────────────────────
\begin{frame}{习题解答 · §2（对偶）}
  \textbf{2.1}~(\diffI)\ $c(i,j)=|i-j|$，$\mu=\nu=\tfrac13\sum_i\delta_i$。最优耦合为恒等：$\pi^\star=\tfrac13\sum_i\delta_{(i,i)}$，成本 $0$。对任意 $(i,i),(j,j)\in\supp\pi^\star$：$c(i,i)+c(j,j)=0\le c(i,j)+c(j,i)=2|i-j|$，故 c-循环单调。

  \medskip
  \textbf{2.2}~(\diffII)\ $\psi(x)=\tfrac12x^2$，$c=\tfrac12|x-y|^2$。$\ctr{\psi}(y)=\inf_x[\tfrac12x^2+\tfrac12(x-y)^2]$；令导数 $2x-y=0\Rightarrow x=y/2$，得 $\ctr{\psi}(y)=y^2/4$。条件 $\ctr{\psi}(y)-\psi(x)=c(x,y)$ 化为 $(x-y/2)^2=0$，故 $\partial_c\psi=\{(x,y):x=y/2\}$（对应映射 $y=2x$）。

  \medskip
  \textbf{2.3}~(\diffIII)\ 由 Thm 5.10(ii)，两个最优 Monge 耦合 $(\Id,T_k)_\#\mu$ 都集中在同一 $\partial_c\psi$。$\mu$ 绝对连续 + 二次成本 $\Rightarrow$ $\partial_c\psi(x)$ 几乎处处单点（Thm 5.30 条件），故 $T_1=T_2$（$\mu$-a.s.）。

  \TLtakeaway{§2：恒等耦合显然 c-循环单调；c-变换是逐点极小化；二次成本 + 绝对连续给出唯一 Brenier 映射。}
\end{frame}

% ── §3 解答 ──────────────────────────────────────────────────────────────────
\begin{frame}{习题解答 · §3（Wasserstein）}
  \textbf{3.1}~(\diffI)\ Bernoulli$(p_0)$、Bernoulli$(q_0)$（$p_0\le q_0$）。分位函数 $F_\mu^{-1}=\mathbf{1}_{t\ge1-p_0}$，$F_\nu^{-1}=\mathbf{1}_{t\ge1-q_0}$，二者之差仅在 $[1-q_0,1-p_0)$ 上为 $1$。故 $\Wp{1}(\mu,\nu)=q_0-p_0$。

  \medskip
  \textbf{3.2}~(\diffII)\ $\mu_k=(1-\tfrac1k)\delta_0+\tfrac1k\delta_{k^2}$。对有界 $\varphi$，$\int\varphi\dd\mu_k\to\varphi(0)$，故 $\mu_k\wkto\delta_0$；但一阶矩 $\int|x|\dd\mu_k=\tfrac1k\cdot k^2=k\to\infty$，违反 Thm 6.9 的矩收敛条件，故 $\Wp{1}(\mu_k,\delta_0)\to\infty$。

  \medskip
  \textbf{3.3}~(\diffIII)\ 代入多维高斯公式 $\Sigma_1=I_n,m_1=0$：$(I^{1/2}\Sigma I^{1/2})^{1/2}=\Sigma^{1/2}$，故 $\Wtwo^2=|m|^2+\mathrm{tr}(I+\Sigma-2\Sigma^{1/2})=|m|^2+\mathrm{tr}((I-\Sigma^{1/2})^2)$。当 $n=1$，$\Sigma=\sigma^2$：$\Wtwo^2=m^2+(1-\sigma)^2$，与正文高斯例吻合。

  \TLtakeaway{§3：一维 Wasserstein 是分位差的 $L^p$ 范数；质量逃逸破坏矩收敛；多维高斯公式在 $n=1$ 退化为标量公式。}
\end{frame}

% ── §4 解答 ──────────────────────────────────────────────────────────────────
\begin{frame}{习题解答 · §4（位移插值）}
  \textbf{4.1}~(\diffI)\ 位移插值 $\mu_t=\delta_t$；线性插值 $(1-t)\delta_0+t\delta_1$。唯一耦合把 $(1-t)$ 的质量从 $0$ 送到 $t$、$t$ 的质量从 $1$ 送到 $t$：$\Wtwo^2=(1-t)t^2+t(1-t)^2=t(1-t)$，故 $\Wtwo=\sqrt{t(1-t)}$（在 $t=\tfrac12$ 取最大 $\tfrac12$）。

  \medskip
  \textbf{4.2}~(\diffII)\ $\mu_0=\mathcal{N}(0,1),\mu_1=\mathcal{N}(0,4)$：(a) $T(x)=(\sigma_1/\sigma_0)x=2x$；(b) $m_t=0$，$\sigma_t=(1-t)\cdot1+t\cdot2=1+t$，即 $\mu_t=\mathcal{N}(0,(1+t)^2)$；(c) $\Wtwo=|\sigma_1-\sigma_0|=1$。

  \medskip
  \textbf{4.3}~(\diffIII)\ Brenier：$T=\nabla\varphi$，高斯下 $\varphi$ 为二次型故 $T(x)=Ax$ 线性、$A$ 对称正定，$A=\Sigma_0^{-1/2}(\Sigma_0^{1/2}\Sigma_1\Sigma_0^{1/2})^{1/2}\Sigma_0^{-1/2}$。(b) $\Sigma_t=((1-t)I+tA)\Sigma_0((1-t)I+tA)^{\top}$。(c) $\Wtwo^2=\mathrm{tr}(\Sigma_0)+\mathrm{tr}(\Sigma_1)-2\,\mathrm{tr}((\Sigma_0^{1/2}\Sigma_1\Sigma_0^{1/2})^{1/2})$。

  \TLtakeaway{§4：线性与位移插值的差距是 $\sqrt{t(1-t)}$；一维高斯方差线性插值；高维高斯由 Bures 几何给出。}
\end{frame}

% ── §5 解答 ──────────────────────────────────────────────────────────────────
\begin{frame}{习题解答 · §5（缩短原理）}
  \textbf{5.1}~(\diffI)\ $x_1=0,x_2=2$，$c=(x-y)^2$。不相交 $(0\to0,2\to2)$ 成本 $=0$；相交 $(0\to2,2\to0)$ 成本 $=4+4=8$；节省 $8$。一般地，差值 $=|x_1-y_2|^2+|x_2-y_1|^2-|x_1-y_1|^2-|x_2-y_2|^2=2(x_1-x_2)(y_1-y_2)\ge0$（当 $x_1<x_2,\,y_1<y_2$ 时两因子同号）。

  \medskip
  \textbf{5.2}~(\diffII)\ $c$-循环单调给出 $(x_1-y_1)^2+(x_2-y_2)^2\le(x_1-y_2)^2+(x_2-y_1)^2$，展开即 $(x_1-x_2)(y_1-y_2)\ge0$。故 $x_1<x_2\Rightarrow x_1-x_2<0\Rightarrow y_1-y_2\le0$，即 $y_1\le y_2$（支撑单调）。

  \medskip
  \textbf{5.3}~(\diffIII)\ (a) $L=|v|^2$ 为 $C^2$ 且 $\nabla_v^2L=2I>0$；$C(\mu_0,\mu_1)<+\infty$ 已设；$\mu_0\in\Pac$ 满足"端点之一绝对连续"。故定理 8.7 适用，$\mu_t\in\Pac$（$0<t<1$）。
  (b) 对称：定理 8.7 对 $\mu_0$ \emph{或} $\mu_1$ 绝对连续均成立。
  (c) 仍成立——改用 $\mu_1\in\Pac$ 那一分支即可，$\mu_0$ 是否原子无关紧要。

  \TLtakeaway{§5：无交叉的代数核是 $(x_1-x_2)(y_1-y_2)\ge0$；定理 8.7 的四条假设（$L\in C^2$、$\nabla_v^2L>0$、成本有限、端点之一 $\in\Pac$）缺一不可。}
\end{frame}

%% ── 技术备用幻灯片 (Backup) ─────────────────────────────
\section*{技术备用幻灯片 · Backup（推迟的技术估计）}

\begin{frame}{（接 §5 定理 8.1）速度 Hölder 估计：设置与双边界}
  \deflab{设置。}
  固定内部时刻 $t_{0}$ 与小窗口 $\tau>0$，记
  $X_{i}=\gamma_{i}(t_{0})$，$V_{i}=\dot{\gamma}_{i}(t_{0})$。Villani 的目标是证明
  $|V_{1}-V_{2}|\le C|X_{1}-X_{2}|^{\beta}$，即式 (8.8)。

  \deflab{上界 (8.21)。}
  把两条曲线在 $[t_{0}-\tau,t_{0}+\tau]$ 的中段互换；\emph{节省}的作用量
  $G:=A[\gamma_{1}]+A[\gamma_{2}]-A[\widetilde{\gamma}_{1}]-A[\widetilde{\gamma}_{2}]\ge0$
  由 Taylor 误差从上方控制：
  \[
    G \;\le\;
    C\!\left(
      \frac{|X_{1}-X_{2}|^{1+\alpha}}{\tau^{1+\alpha}}
      +\tau |V_{1}-V_{2}|^{1+\alpha}
    \right).
  \]

  \deflab{下界 (8.22)。}
  假设 (iii) 的 $(2+\kappa)$-凸性给出\emph{同一个} $G$ 的下界：
  \[
    G \;\ge\;
    K\bigl(|V_{1}-V_{2}|-A\tau |X_{1}-X_{2}|\bigr)^{2+\kappa}.
  \]

  \TLtakeaway{上界来自 shortcut 的 Taylor 误差；下界来自速度方向的严格凸性。两者夹住同一个节省量 $G\ge0$。}
\end{frame}

\begin{frame}{（接 §5 定理 8.1）速度 Hölder 估计：比较与优化}
  \deflab{情形一。}
  若 $|V_{1}-V_{2}|\le 2A\tau |X_{1}-X_{2}|$，则速度差已被位置差控制；后面选择允许的 $\tau$ 后即可收尾。

  \deflab{情形二。}
  否则有 $|V_{1}-V_{2}|-A\tau |X_{1}-X_{2}|\ge |V_{1}-V_{2}|/2$。比较上一帧的上下界，得到 (8.23)：
  \[
    |V_{1}-V_{2}|^{2+\kappa}
    \le
    C\!\left(
      \frac{|X_{1}-X_{2}|^{1+\alpha}}{\tau^{1+\alpha}}
      +\tau |V_{1}-V_{2}|^{1+\alpha}
    \right).
  \]

  \deflab{优化。}
  当 $|V_{1}-V_{2}|>0$ 时，取 $\tau\sim|V_{1}-V_{2}|^{1+\kappa-\alpha}$
  以吸收右端第二项（Villani 第 178 页），得到
  \[
    |V_{1}-V_{2}|\le C|X_{1}-X_{2}|^{\beta},
    \qquad \beta=\beta(\alpha,\kappa)>0.
  \]
  在二次成本 $c=d^{2}$ 中，$\alpha=1,\kappa=0$，于是 $\beta=1$。

  \TLtakeaway{这是 §5 "定理 8.1 证明思路（凸性矛盾）" 暂缓的定量估计：速度差由中间位置差的幂控制。}
\end{frame}

\begin{frame}{（接 §2 定理 5.10）可测性补证：势函数为何可测}
  \deflab{问题。}
  链式公式 (5.22) 中的 $\psi$ 是不可数个 usc 函数的上确界，不自动是 Borel 函数。

  \deflab{Egorov + 内正则。}
  取 $0\le c_{\ell}\nearrow c$ 为连续逼近。对每个 $k$，可取紧集 $\Gamma_{k}\subset\Gamma$，使得
  \[
    \pi[\Gamma_{k}]\ge 1-\frac{2}{k},\qquad
    c_{\ell}\to c\ \mathrm{uniformly\ on}\ \Gamma_{k},\qquad
    \Gamma=\bigcup_{k}\Gamma_{k}.
  \]

  \deflab{截断势函数。}
  链长截为 $m$，链点限制在 $\Gamma_{k}$，横跳项用 $c_{\ell}$，得 $\psi_{m,k,\ell}(x)$。
  它关于 $x$ 为 lsc；逐点极限是 Borel，且
  \[
    \psi(x)=\sup_{m\in\N}\sup_{k\in\N}\lim_{\ell\to\infty}\psi_{m,k,\ell}(x).
  \]
  右端只含可数次上确界与极限，故 $\psi$ 可测。

  \TLtakeaway{Egorov + 内正则把不可数上确界转化为可数紧集族上的截断链极限，从而得到 $\psi$ 可测。}
\end{frame}

\begin{frame}{（接 §2 定理 5.10）可测性补证：$\phi$ 与可测选择}
  \deflab{$\phi=\ctr{\psi}$ 的可测版本。}
  分解 $\pi$ 为以第二边际为基的条件测度：
  \[
    \pi(\dd x\,\dd y)=\pi(\dd x\mid y)\,\nu(\dd y),
    \qquad
    \widetilde{\phi}(y)=\int_{\X}[\psi(x)+c(x,y)]\,\pi(\dd x\mid y).
  \]
  则 $\phi(y)=\widetilde{\phi}(y)$ 对 $\nu$-a.e.\ 的 $y$ 成立。

  \deflab{$\partial_{c}\psi$ 的可测版本。}
  除去一个 $\pi$-零集后，$\partial_{c}\psi$ 与下列 Borel 集一致：
  \[
    \{(x,y):\widetilde{\phi}(y)-\psi(x)=c(x,y)\}.
  \]

  \deflab{推论 5.22（最优计划的可测选择）。}
  若 $c$ 连续、$\inf c>-\infty$，且 $\omega\mapsto(\mu_{\omega},\nu_{\omega})$ 可测，则可选可测的最优计划 $\omega\mapsto\pi_{\omega}$。
  证明思路：最优计划全体 $O$ 是 Polish；边际映射 $\Phi:O\to\Pp(\X)\times\Pp(\Y)$ 连续且 onto，纤维紧，故可用可测选择定理。

  \TLtakeaway{这补上 §2 "定理 5.10 证明思路（二c）" 暂缓的可测性缺口：势函数、c-次微分与最优计划选择都可取为可测对象。}
\end{frame}
